% Part: lambda-calculus
% Chapter: syntax
% Section: term-revisited

\documentclass[../../../include/open-logic-section]{subfiles}

\begin{document}

\olfileid{lam}{syn}{tr}
\olsection{الحدود بوصفها أصناف تكافؤ بحسب~$\alpha$}

من الآن ندرس الحدود مع إهمال الفرق الذي لا يتجاوز التكافؤ بحسب~$\alpha$.
فإذا كتبنا حدًا قصدنا صنف تكافؤه بحسب~$\alpha$.
فكتابتنا~$\lambd[a][\lambd[b][a c]]$ تعني مجموعة الحدود التي تكافئه بحسب~$\alpha$،
ومنها~$\lambd[a][\lambd[b][a c]]$ و$\lambd[b][\lambd[a][b c]]$ وما شابهها.

وكانت حروف مثل~$N, Q$ في الأقسام الماضية أسماء لحدود مفردة؛
أما فيما يأتي فنسمي بها الأصناف، وهي التي تكون موضوع البحث بدل الحدود المفردة.
وتبقى حروف مثل~$x, y$ أسماء للمتغيرات.

ونجعل~$\rep{M}$ رمزًا لـ!!{element} كيفما اتفق من الصنف~$M$.
فإن احتجنا إلى ممثلين متعددين استعملنا~$\rep{M}[0], \rep{M}[1]$، وهكذا.

ولإيجاز العبارة نبقي ترميز الحدود، ونعرّف معناه على الأصناف:
\begin{defn}
  \begin{enumerate}
  \item $\lambd[x][N]$ هو، بالتعريف، الصنف الذي ينتمي إليه
    $\lambd[x][\rep{N}]$.
  \item $PQ$ هو، بالتعريف، الصنف الذي ينتمي إليه~$\rep{P}\rep{Q}$.
  \end{enumerate}
\end{defn}

وهذان التعريفان لا يتوقفان على اختيار الممثل؛ فإن تحويل~$\alpha$ متوافق، ومن ذلك يظهر حسن تعريفهما.

\begin{defn} \ollabel{def:fv}
  مجموعة \emph{المتغيرات الحرة} لصنف التكافؤ بحسب~$\alpha$ المسمى~$M$،
  ورمزها~$\FV{M}$، هي بالتعريف~$\FV{\rep{M}}$.
\end{defn}

ولا تتعلق هذه المجموعة باختيار الممثل، لأن~$\FV{\rep{M}[0]} =
\FV{\rep{M}[1]}$، كما أثبتنا في~\olref[alp]{thm:fv}.

وننقل ترميز الإحلال إلى الأصناف كذلك:
\begin{defn} \ollabel{def:sub}
  \emph{إحلال}~$R$ محل~$y$ في~$M$، ورمزه~$\Subst{M}{R}{y}$،
  هو صنف التكافؤ بحسب~$\alpha$ الذي ينتمي إليه
  $\Subst{\rep{M}}{\rep{R}}{y}$؛ ونأخذ~$\rep{M}$ و$\rep{R}$ أي ممثلين يكون الإحلال عليهما معرّفًا.
\end{defn}

وحسن تعريف هذا الإحلال ثابت أيضًا بـ\olref[alp]{cor:sub}.

ويتضح مقدار ما ييسره هذا التعريف من الاستدلال في المثال الآتي:
\begin{align}
  \Subst{\lambd[x][x]}{y}{x} & \ollabel{eq:1}\\
  &= \Subst{\lambd[z][z]}{y}{x} \ollabel{eq:2}\\
  &= \lambd[z][\Subst{z}{y}{x}] \\
  &= \lambd[z][z]
\end{align}

لو حملنا~\olref{eq:1} على الإحلال في الحدود المفردة لكان غير معرّف.
لكن المراد الآن الإحلال في الأصناف، كما تقدم؛ ولذلك جاز الانتقال إلى~\olref{eq:2}:
فلنا أن نستبدل~$\lambd[x][x]$ بـ$\lambd[z][z]$، إذ كلاهما في صنف واحد.

وللعلة نفسها نفترض فيما يأتي أن اختيار الممثلين قد استوفى شروط الإحلال.
فإذا ورد~$\Subst{\lambd[x][N]}{R}{y}$، مثلًا، حملناه على أن الممثل~$\lambd[x][N]$
اختير بحيث~$x \neq y$ و$x \notin \FV{R}$.

ولما كان إطلاق «صنف» على~$\lambd[x][x]$ غير مألوف بعض الشيء،
سمينا هذه الأشياء فيما يأتي حدود~$\Lambda$، أو «حدودًا» بإطلاق في بقية هذا الجزء،
تمييزًا لها من حدود~$\lambda$ المعهودة.

\begin{editorial}
  لا نستغني بذلك عن الحدود الأصلية؛ فتعريف حدود~$\Lambda$ كله يرجع إلى
  حدود~$\lambda$. ولم نضع بعد طريقة لتعريف الدوال على حدود~$\Lambda$،
  فلا بد من تعريف كل دالة أولًا على حدود~$\lambda$،
  ثم إنزالها إلى حدود~$\Lambda$، على نحو ما صنعنا بالإحلال.
  غير أننا نفترض أن القارئ يدرك حدسيًا كيف تعرّف الدوال على حدود~$\Lambda$.
\end{editorial}
\end{document}
